% part: intuitionistic-logic
% chapter: introduction
% section: natural-deduction

\documentclass[../../../include/open-logic-section]{subfiles}

\begin{document}

\olfileid{int}{int}{ntd}

\olsection{استنتاج طبیعی}

استنتاج طبیعی بدون قواعد~$\FalseCl$ یک دستگاه استاندارد
!!{derivation} برای منطق شهودگرایانه است. قواعد را در اینجا تکرار
می‌کنیم و انگیزهٔ آن‌ها را با تفسیر BHK نشان می‌دهیم. در هر مورد، می‌توان
قاعده را چنین فهمید که به ما اجازه می‌دهد نتیجه بگیریم اگر مقدمات
ساخت‌هایی داشته باشند، نتیجه نیز ساختی دارد.

چون !!{derivation}های استنتاج طبیعی فرض‌های باز دارند،
باید چنین !!a{derivation}، مثلاً اشتقاق~$!A$ از فرض‌های
!!{undischarged}~$\Gamma$، را تابعی بدانیم که ساخت‌های همهٔ
$!B \in \Gamma$ را به ساختی از~$!A$ تبدیل می‌کند. اگر
!!a{derivation}ای از~$!A$ بدون هیچ فرض !!{undischarged} وجود داشته
باشد، آنگاه ساختی از~$!A$ به معنای تفسیر BHK وجود دارد. بااین‌حال، برای
بحث حاضر هرگاه نیازی نباشد~$\Gamma$ را نمی‌نویسیم.

یک فرض~$!A$ به‌تنهایی !!a{derivation}ای از~$!A$ از فرض
!!{undischarged}~$!A$ است. این با تفسیر BHK سازگار است: تابع همانی روی
ساخت‌ها هر ساختِ~$!A$ را به ساختی از~$!A$ تبدیل می‌کند.

\subsection{عطف}

\begin{defish}
\AxiomC{$!A$}
\AxiomC{$!B$}
\RightLabel{\Intro{\land}}
\BinaryInfC{$!A \land !B$}
\DisplayProof
\hfill
\begin{tabular}{l}
\AxiomC{$!A \land !B$}
\RightLabel{\Elim{\land}}
\UnaryInfC{$!A$}
\DisplayProof
\\[3ex]
\AxiomC{$!A \land !B$}
\RightLabel{\Elim{\land}}
\UnaryInfC{$!B$}
\DisplayProof
\end{tabular}
\end{defish}

فرض کنید به‌ترتیب ساخت‌های~$N_1$ و~$N_2$ از~$!A_1$ و~$!A_2$ را داریم.
آنگاه ساختی از~$!A_1 \land !A_2$ نیز داریم، یعنی جفت
$\tuple{N_1, N_2}$.

در تفسیر BHK، ساختِ~$!A_1 \land !A_2$ جفت~$\tuple{N_1, N_2}$ است.
پس فرض کنید چنین جفتی داریم. آنگاه ساختی از هریک از دو طرف عطف نیز
داریم: $N_1$ ساختی از~$!A_1$ و~$N_2$ ساختی از~$!A_2$ است.

\subsection{شرطی}

\begin{defish}
  \AxiomC{$\Discharge{!A}{u}$}
  \DeduceC{$!B$}
  \DischargeRule{$\Intro{\lif}$}{u}
  \UnaryInfC{$!A \lif !B$}
  \DisplayProof
\hfill
  \AxiomC{$!A \lif !B$}
  \AxiomC{$!A$}
  \RightLabel{$\Elim{\lif}$}
  \BinaryInfC{$!B$}
  \DisplayProof
\end{defish}

اگر !!a{derivation}ای از~$!B$ از فرض !!{undischarged}~$!A$ داشته
باشیم، تابع~$f$ای وجود دارد که ساخت‌های~$!A$ را به ساخت‌های~$!B$ تبدیل
می‌کند. همان تابع ساختی از~$!A \lif !B$ است. پس اگر مقدمهٔ~$\Intro\lif$
مشروط به ساختی از~$!A$ ساختی داشته باشد، نتیجهٔ~$!A \lif !B$ نیز
ساختی دارد.

از سوی دیگر، فرض کنید ساخت~$N$ از~$!A$ و ساخت~$f$ از~$!A \lif !B$
وجود داشته باشد. ساختِ~$!A \lif !B$ تابعی است که ساخت‌های~$!A$ را به
ساخت‌های~$!B$ تبدیل می‌کند. پس~$f(N)$ ساختی از~$!B$ است؛ یعنی نتیجهٔ
$\Elim\lif$ ساختی دارد.

\subsection{فصل}

\begin{defish}
\begin{tabular}{l}
\AxiomC{$!A$}
\RightLabel{\Intro{\lor}}
\UnaryInfC{$!A \lor !B$}
\DisplayProof
\\[3ex]
\AxiomC{$!B$}
\RightLabel{\Intro{\lor}}
\UnaryInfC{$!A \lor !B$}
\DisplayProof
\end{tabular}
\hfill
\AxiomC{$!A \lor !B$}
\AxiomC{$\Discharge{!A}{n}$}
\DeduceC{$!C$}
\AxiomC{$\Discharge{!B}{n}$}
\DeduceC{$!C$}
\DischargeRule{\Elim{\lor}}{n}
\TrinaryInfC{$!C$}
\DisplayProof
\end{defish}

اگر ساخت~$N_i$ از~$!A_i$ را داشته باشیم، می‌توانیم آن را به ساخت
$\tuple{i, N_i}$ از~$!A_1 \lor !A_2$ تبدیل کنیم. از سوی دیگر، فرض کنید
ساختی از~$!A_1 \lor !A_2$، یعنی جفت~$\tuple{i, N_i}$ که در آن~$N_i$
ساختی از~$!A_i$ است، و نیز توابع~$f_1$ و~$f_2$ را داشته باشیم که
ساخت‌های~$!A_1$ و~$!A_2$ را به‌ترتیب به ساخت‌های~$!C$ تبدیل کنند.
آنگاه~$f_i(N_i)$ ساختی از~$!C$، یعنی نتیجهٔ~$\Elim\lor$، است.

\subsection{محال}

\begin{defish}
  \AxiomC{$\lfalse$}
  \RightLabel{\FalseInt}
  \UnaryInfC{$!A$}
  \DisplayProof
\end{defish}

اگر !!a{derivation}ای از~$\lfalse$ از فرض‌های !!{undischarged}
$!B_1$، \dots، $!B_n$ داشته باشیم، تابع~$f(M_1, \dots, M_n)$ای وجود
دارد که ساخت‌های~$!B_1$، \dots، $!B_n$ را به ساختی از~$\lfalse$ تبدیل
می‌کند. چون~$\lfalse$ هیچ ساختی ندارد، نمی‌تواند ساختی برای همهٔ
$!B_1$، \dots، $!B_n$ نیز وجود داشته باشد. ازاین‌رو، $f$ این خاصیت را
نیز دارد که \emph{اگر} $M_1$، \dots، $M_n$ به‌ترتیب ساخت‌های~$!B_1$،
\dots، $!B_n$ باشند، \emph{آنگاه} $f(M_1, \dots, M_n)$ ساختی از~$!A$
است.

\subsection{قواعد~$\lnot$}

از آنجا که~$\lnot !A$ به‌صورت~$!A \lif \lfalse$ تعریف شده است، به‌دقت
سخن بگوییم به قواعدی برای~$\lnot$ نیاز نداریم. اما اگر چنین قواعدی
می‌داشتیم، صورتشان چنین بود:

\begin{defish}
\AxiomC{$\Discharge{!A}{n}$}
\noLine
\DeduceC{$\lfalse$}
\DischargeRule{\Intro{\lnot}}{n}
\UnaryInfC{$\lnot !A$}
\DisplayProof
\hfill
\AxiomC{$\lnot !A$}
\AxiomC{$!A$}
\RightLabel{\Elim{\lnot}}
\BinaryInfC{$\lfalse$}
\DisplayProof
\end{defish}


\subsection{نمونه‌های \usetoken{P}{derivation}}

\begin{enumerate}
\item $\Proves !A \lif (\lnot !A \lif \lfalse)$، یعنی
  $\Proves !A \lif ((!A \lif \lfalse) \lif \lfalse)$
  \begin{prooftree}
    \AxiomC{$\Discharge{!A}{2}$}
    \AxiomC{$\Discharge{!A \lif \lfalse}{1}$}
    \RightLabel{\Elim\lif}
    \BinaryInfC{$\lfalse$}
    \DischargeRule{\Intro\lif}{1}
    \UnaryInfC{$(!A \lif \lfalse)\lif \lfalse$}
    \DischargeRule{\Intro\lif}{2}
    \UnaryInfC{$!A \lif (!A \lif \lfalse) \lif \lfalse$}
  \end{prooftree}

\item $\Proves ((!A \land !B) \lif !C) \lif (!A \lif (!B \lif !C))$
  \begin{prooftree}
    \AxiomC{$\Discharge{(!A \land !B) \lif !C}{3}$}
    \AxiomC{$\Discharge{!A}{2}$}
    \AxiomC{$\Discharge{!B}{1}$}
    \RightLabel{\Intro\land}
    \BinaryInfC{$!A \land !B$}
    \RightLabel{\Elim\lif}
    \BinaryInfC{$!C$}
    \DischargeRule{\Intro\lif}{1}
    \UnaryInfC{$!B \lif !C$}
    \DischargeRule{\Intro\lif}{2}
    \UnaryInfC{$!A \lif (!B \lif !C)$}
    \DischargeRule{\Intro\lif}{3}
    \UnaryInfC{$((!A \land !B) \lif !C) \lif (!A \lif (!B \lif !C))$}
  \end{prooftree}

\item $\Proves \lnot(!A \land \lnot !A)$، یعنی
  $\Proves (!A \land (!A \lif \lfalse))\lif \lfalse$
  \begin{prooftree}
    \AxiomC{$\Discharge{!A \land (!A \lif \lfalse)}{1}$}
    \RightLabel{\Elim\land}
    \UnaryInfC{$!A \lif \lfalse$}
    \AxiomC{$\Discharge{!A \land (!A \lif \lfalse)}{1}$}
    \RightLabel{\Elim\land}
    \UnaryInfC{$!A$}
    \RightLabel{\Elim\lif}
    \BinaryInfC{$\lfalse$}
    \DischargeRule{\Intro\lif}{1}
    \UnaryInfC{$(!A \land (!A \lif \lfalse))\lif \lfalse$}
  \end{prooftree}

\item $\Proves \lnot\lnot(!A \lor \lnot !A)$، یعنی
  $\Proves ((!A \lor (!A \lif \lfalse))\lif \lfalse)\lif \lfalse$
  \begin{prooftree}\hskip-3em
    \AxiomC{$\Discharge{(!A \lor (!A \lif \lfalse))\lif \lfalse}{2}$}
    \AxiomC{$\Discharge{(!A \lor (!A \lif \lfalse))\lif \lfalse}{2}$}
    \AxiomC{$\Discharge{!A}{1}$}
    \RightLabel{\Intro\lor}
    \UnaryInfC{$!A \lor (!A \lif \lfalse)$}
    \RightLabel{\Elim\lif}
    \BinaryInfC{$\lfalse$}
    \DischargeRule{\Intro\lif}{1}
    \UnaryInfC{$!A \lif \lfalse$}
    \RightLabel{\Intro\lor}
    \UnaryInfC{$!A \lor (!A \lif \lfalse)$}
    \RightLabel{\Elim\lif}
    \insertBetweenHyps{\hskip -4em}
    \BinaryInfC{$\lfalse$}
    \DischargeRule{\Intro\lif}{2}
    \UnaryInfC{$((!A \lor (!A \lif \lfalse))\lif \lfalse)\lif \lfalse$}
  \end{prooftree}
\end{enumerate}

\begin{prop}
  اگر~$\Gamma \Proves !A$ در استنتاج طبیعی برای منطق شهودگرایانه،
  آنگاه~$\Gamma \Proves !A$ در منطق کلاسیک نیز. به‌ویژه، اگر~$!A$
  قضیه‌ای شهودگرایانه باشد، قضیه‌ای کلاسیک نیز هست.
\end{prop}

\begin{proof}
  هر قاعدهٔ استنتاج طبیعی قاعده‌ای در استنتاج طبیعی کلاسیک نیز هست؛ پس
  هر !!{derivation} در منطق شهودگرایانه !!a{derivation}ای در منطق
  کلاسیک نیز هست.
\end{proof}

\begin{prob}
  در منطق شهودگرایانه، !!{derivation}هایی از !!{formula}های زیر بدهید:
  \begin{enumerate}
    \item $(\lnot !A \lor !B) \lif (!A \lif !B)$
    \item $\lnot\lnot\lnot !A \lif \lnot !A$
    \item $\lnot\lnot (!A \land !B) \liff (\lnot\lnot !A\land \lnot \lnot !B)$
    \item $\lnot (!A \lor !B) \liff (\lnot !A \land \lnot !B)$
    \item $(\lnot !A \lor \lnot !B) \lif \lnot (!A \land !B)$
    \item $\lnot\lnot ( !A \land !B) \lif  (\lnot\lnot !A \lor
    \lnot\lnot !B)$
  \end{enumerate}
\end{prob}

\end{document}
